\documentclass[11pt]{article}
\usepackage{amsmath,amssymb,amsthm,hyperref}
\newtheorem{theorem}{Theorem}
\newtheorem{lemma}[theorem]{Lemma}
\newcommand{\E}{\mathbb E}
\newcommand{\Prob}{\mathbb P}
\title{A refined uniform bound for random palindrome blocks}
\author{Research note; independently reviewed by two other agents}
\date{30 September 2026}
\begin{document}
\maketitle

Concatenate $m$ independently and uniformly relabelled palindrome blocks
$\rho_i\rho_i^{\mathrm{rev}}$ on $n$ letters. Each resulting die has $2m$
faces. Write $R_\sigma=n!P(\sigma)$ for an output permutation $\sigma$.

\begin{theorem}\label{thm:main}
For $n\geq3$ and $0<\varepsilon\leq1$, put
\[
 p=\max\left(2,\frac{\log(2n!)}{\log3}\right).
\]
If
\[
 m\geq 8n^{11/10}p^{1/5}\varepsilon^{-2/5},
\]
then with probability at least $1/2$,
$|R_\sigma-1|\leq\varepsilon$ simultaneously for all $\sigma\in S_n$.
Consequently
$O(n^{13/10}(\log n)^{1/5}\varepsilon^{-2/5})$
faces per die suffice for pointwise relative approximate fairness.
\end{theorem}

All logarithms are natural. There are no asymptotic restrictions on $n$
or $\varepsilon$ in the theorem. The proof controls the entire nonlinear
expansion, including segments of length greater than three.

\section{Two elementary estimates}
\begin{lemma}[Bounded multilinear chaos]\label{lem:chaos}
Suppose independent centered random variables $Z_v$ satisfy $|Z_v|\leq b_v$.
For a multilinear homogeneous polynomial
$T=\sum_{|A|=\ell}c_A\prod_{v\in A}Z_v$ and $p\geq2$,
\[
 \|T\|_p\leq(2\sqrt{p-1})^\ell
      \left(\sum_{|A|=\ell}c_A^2\prod_{v\in A}b_v^2\right)^{1/2}.
\]
\end{lemma}
\begin{proof}
Introduce independent copies $Z_v'$ and replace each $Z_v$ by $Z_v-Z_v'$.
Conditional expectation over the copies recovers $T$, so Jensen applies.
Each independent pair $(Z_v,Z_v')$ can be swapped, which multiplies its
difference by an independent Rademacher sign and preserves the joint law.
Conditional on the pairs, the degree-$\ell$ Rademacher
hypercontractive inequality \cite{bonami,odonnell} bounds the $L^p$ norm by
$(p-1)^{\ell/2}$ times the coefficient $\ell^2$ norm. Every difference is
bounded by $2b_v$. The resulting bound is uniform in the pairs.
\end{proof}

For $g=(g_0,\ldots,g_\ell)$ a weak composition of $M>0$, let
\[
 C_d(g/M)=\sum_{a_0+\cdots+a_\ell=d}
 \left[\frac{d!}{a_0!\cdots a_\ell!}
       \prod_{j=0}^\ell(g_j/M)^{a_j}\right]^2
\]
be the collision probability of a multinomial distribution; $0^0=1$.
\begin{lemma}[A lattice collision bound]\label{lem:lattice}
For integers $\ell\geq1$, $d\geq0$ and $M\geq d+1$,
\[
 \sum_{g_0+\cdots+g_\ell=M}C_d(g/M)
 \leq\left(\frac{108M}{\sqrt{d+1}}\right)^\ell.
\]
\end{lemma}
\begin{proof}
The maximum mass of a Poisson distribution of mean $\lambda\geq0$ is
at most $2/\sqrt{1+\lambda}$, and the probability of $d$ under
$\operatorname{Poisson}(d)$ is at least $1/(3\sqrt{d+1})$.
These follow directly from Stirling's elementary factorial inequalities
(the cases $d=0$ and $\lambda\leq1$ are immediate).
Condition independent Poisson variables of means $dp_j$ on their sum
being $d$. Since a collision probability is at most the maximum mass,
\[
 C_d(p)\leq3\,2^{\ell+1}\sqrt{d+1}
              \prod_{j=0}^\ell(1+dp_j)^{-1/2}.
\]
For every $g$, some coordinate satisfies $g_j\geq M/(\ell+1)$.
For this coordinate,
$\sqrt{d+1}(1+dg_j/M)^{-1/2}\leq\sqrt{\ell+1}$.
Sum over the possible choices of $j$, then drop the sum constraint on the
remaining $\ell$ coordinates, each of which lies between $0$ and $M$.
For the decreasing function $w(x)=(1+dx/M)^{-1/2}$,
\[
 \sum_{g=0}^M w(g)\leq 1+\int_0^M w(x)\,dx
 \leq 1+\frac{2M}{\sqrt{d+1}}
 \leq\frac{3M}{\sqrt{d+1}}.
\]
The middle bound also holds for $d=0$, and the last uses
$M\geq d+1\geq\sqrt{d+1}$. Hence the desired sum is at most
\[
 6(\ell+1)^{3/2}
     \left(\frac{6M}{\sqrt{d+1}}\right)^\ell
 \leq\left(\frac{108M}{\sqrt{d+1}}\right)^\ell,
\]
because $\ell+1\leq2^\ell$ and $6\leq6^\ell$.
\end{proof}

\section{The exact expansion and a coloring argument}
Fix a target order $\sigma$. Let $L$ have the multinomial law with $n$ trials
and $m$ equally likely cells. Split $\sigma$ into consecutive segments $A_i$
of lengths $L_i$. If $q_i(A)$ is the probability of the prescribed order
inside block $i$, set
\[
 \delta_i(A)=|A|!q_i(A)-1.
\]
For empty segments set $\delta_i=0$. Then, exactly,
\[
 R_\sigma=\E_L\prod_i(1+\delta_i(A_i)).
\]
Uniform relabelling gives $\E\delta_i(A)=0$. For $|A|\leq2$ the value is
identically zero. For $k=|A|\geq3$,
\[
 |\delta_i(A)|\leq b_k:=k!2^{1-k}.
\]
Indeed any prescribed order occurs at most twice as a subsequence of one
palindrome, so $q_i(A)\leq2^{1-k}$; also $b_k\geq1$.

Represent $\rho_i$ by the relative order of independent continuous
priorities $Y_{i,1},\ldots,Y_{i,n}$, indexed by the positions of $\sigma$.
Write $\delta_{i,a,k}$ for the segment starting at position $a$ and having
length $k$. For fixed $k$, the variables $\delta_{i,a,k}$ with $a$ in a
fixed residue class modulo $k$ are independent: they use disjoint sets of
priorities. Blocks are independent of each other as well.

The degree-$\ell$ part of $R_\sigma-1$ has the form
\[
 T_\ell=\sum_{I,\,\boldsymbol k,\,\boldsymbol a}
       c(I,\boldsymbol k,\boldsymbol a)
       \prod_{j=1}^\ell\delta_{i_j,a_j,k_j},\qquad
 I=\{i_1<\cdots<i_\ell\},\quad k_j\geq3.
\]
The coefficients are nonnegative probabilities of the corresponding block
occupancy and segment start events. The starts must specify disjoint
segments in their natural order.

Here is a useful bound for any restriction $\mathcal D$ on the indices of
this sum:
\begin{equation}\label{eq:colored}
 \|T_\ell^{\mathcal D}\|_p\leq(2\sqrt{p-1})^\ell
 \left[\sum_{(I,\boldsymbol k,\boldsymbol a)\in\mathcal D}
 c(I,\boldsymbol k,\boldsymbol a)^2
          \prod_{j=1}^\ell\frac{k_j b_{k_j}^2}{q_{k_j}}\right]^{1/2},
 \qquad q_k=2^{-(k-2)},\quad k\geq3.
\end{equation}
To prove it, independently for each block choose a size $K_i$ with
probabilities $q_k$, and a uniform residue $B_i$ modulo $K_i$. Keep only
terms with $k_j=K_{i_j}$ and $a_j\equiv B_{i_j}\pmod{k_j}$, multiplying
their coefficients by $\prod_jk_j/q_{k_j}$.
Averaging over these auxiliary choices recovers the original sum.
Given the choices, all remaining variables are independent centered bounded
random variables, so Lemma~\ref{lem:chaos} applies.
Minkowski followed by Jensen for the square root gives \eqref{eq:colored},
since the probability of retaining a term is $\prod_jq_{k_j}/k_j$.
Sizes $K_i>n$ simply retain no terms.

For fixed $I$ and lengths $\boldsymbol k$, write $s=\sum_jk_j$, $d=n-s$,
and let
\[
 g_0=i_1-1,\quad g_j=i_{j+1}-i_j-1\ (1\leq j<\ell),\quad
 g_\ell=m-i_\ell,\qquad M=\sum_jg_j=m-\ell.
\]
The unspecified blocks fall into these $\ell+1$ gaps. A choice of segment
starts is equivalent to counts $a_0,\ldots,a_\ell$ of remaining labels in
these gaps, summing to $d$. Direct summation of the multinomial law gives
\begin{equation}\label{eq:coefficient}
 c(I,\boldsymbol k,\boldsymbol a)
 =\frac{n^{\underline{s}}}{m^s\prod_jk_j!}
       (M/m)^d
       \frac{d!}{a_0!\cdots a_\ell!}
       \prod_{j=0}^\ell(g_j/M)^{a_j}.
\end{equation}
We henceforth use $\boldsymbol a$ for these gap counts rather than the
equivalent starts. The map $I\mapsto\boldsymbol g$ is a bijection between
$\ell$-subsets of $[m]$ and weak compositions of $m-\ell$.

\section{Short total support}
Split the expansion according to $s\leq n/2$ and $s>n/2$, writing their
sums as $T^{\mathrm{short}}$ and $T^{\mathrm{long}}$.
Assume $m\geq n$, as guaranteed in the theorem. For short support,
$d\geq n/2$ and
$M=m-\ell\geq d+2\ell\geq d+1$.
Using \eqref{eq:coefficient}, the square sum in \eqref{eq:colored} is at most
\[
 \sum_{\boldsymbol k:\,s\leq n/2}
 \prod_{j=1}^\ell\left[
       (n/m)^{2k_j}2^{2-2k_j}\frac{k_j}{q_{k_j}}\right]
 \sum_{g_0+\cdots+g_\ell=M}C_d(g/M).
\]
Here we dropped $(M/m)^{2d}\leq1$ and used
$n^{\underline{s}}\leq n^s$. Lemma~\ref{lem:lattice} bounds the collision
sum by $(160m/\sqrt n)^\ell$, since $108\sqrt2<160$.
The length sum separates after extending it to all $k_j\geq3$. For
$z=n^2/(2m^2)\leq1/2$,
\[
 \sum_{k\geq3}(n/m)^{2k}2^{2-2k}\frac{k}{q_k}
 =\sum_{k\geq3}kz^k
 =\frac{z^3(3-2z)}{(1-z)^2}
 \leq8z^3=\frac{n^6}{m^6}.
\]
Consequently, putting
\[
 a=2\sqrt{160}\,\frac{n^{11/4}\sqrt p}{m^{5/2}},
\]
we have
\begin{equation}\label{eq:short}
 \|T^{\mathrm{short}}\|_p\leq\sum_{\ell\geq1}a^\ell
 =\frac a{1-a}\qquad(a<1).
\end{equation}

\section{Long total support}
For long support, bound each collision probability by one and sum over
$I$, obtaining $\binom m\ell\leq m^\ell/\ell!$.
For any $\theta\geq1$, the condition $s>n/2$ allows us to multiply the
nonnegative summand by the upper bound $\theta^{s-n/2}$.
Equation~\eqref{eq:colored} gives
\[
 \|T_\ell^{\mathrm{long}}\|_p
 \leq\theta^{-n/4}\frac{A_\theta^\ell}{\sqrt{\ell!}},
 \qquad
 A_\theta=2\sqrt{pm}
   \left(\sum_{k\geq3}k
           \left(\frac{\theta n^2}{2m^2}\right)^k\right)^{1/2}.
\]
Set $t=m/n\geq1$ and $\theta=t^{1/3}$. The power-series argument is
$1/(2t^{5/3})\leq1/2$, so
\[
 A_\theta^2\leq4pm\,t^{-5}=\frac{4pn}{t^4}.
\]
Cauchy--Schwarz gives
$\sum_{\ell\geq0}A^\ell/\sqrt{\ell!}\leq\sqrt2e^{A^2}$.
Thus
\begin{equation}\label{eq:long}
 \|T^{\mathrm{long}}\|_p
 \leq\sqrt2\exp\left(-\frac n{12}\log t+\frac{4pn}{t^4}\right).
\end{equation}

\section{Constants and the simultaneous guarantee}
The theorem's assumption implies
$t\geq8n^{1/10}p^{1/5}\varepsilon^{-2/5}\geq8$ and $a\leq\varepsilon/6$,
because $2\sqrt{160}/8^{5/2}<1/6$.
Equation~\eqref{eq:short} therefore gives
$\|T^{\mathrm{short}}\|_p\leq\varepsilon/5$.
Also $p\leq n^2$ for $n\geq3$, whence
\[
 \frac{4p}{t^4}\leq\frac1{1024}(p/n^2)^{1/5}
                  \varepsilon^{8/5}
 \leq\frac1{1024}\leq\frac{\log t}{24}.
\]
For $n\geq60$, \eqref{eq:long} now yields
\[
 \|T^{\mathrm{long}}\|_p
 \leq\sqrt2\,t^{-n/24}
 \leq\sqrt2\,8^{-5/2}\varepsilon
 =\frac\varepsilon{128}.
\]
We conclude that $\|R_\sigma-1\|_p<\varepsilon/3$.

For completeness the finitely many cases $3\leq n<60$ follow from the
simpler blockwise canonical-kernel bound, reproduced here. If $F_I$ is the
part of the product expansion using exactly the selected blocks $I$, then
its separate centering and the selected-cell multinomial marginal give
\[
 \|F_I\|_\infty\leq\beta^{|I|},\qquad
 \beta=\sum_{k\geq3}2^{1-k}(n/m)^k
      =\frac{n^3}{4m^3(1-n/(2m))}\leq\frac{n^3}{2m^3}.
\]
Iterated symmetrization in the independent block variables, followed by
the same Rademacher hypercontractivity argument, gives
\[
 \|R_\sigma-1\|_p\leq
 \sum_{\ell\geq1}\frac{b^\ell}{\sqrt{\ell!}}
 \leq\frac b{1-b},\qquad
 b=2\beta\sqrt{mp}\leq\frac{n^3\sqrt p}{m^{5/2}}
 \leq\frac{60^{1/4}}{8^{5/2}}\varepsilon<\frac\varepsilon4.
\]
Thus again $\|R_\sigma-1\|_p\leq\varepsilon/3$.
For every $n\geq3$, Markov and a union bound give
\[
 \Prob\bigl(\exists\sigma:\ |R_\sigma-1|>\varepsilon\bigr)
 \leq n!3^{-p}\leq\frac12.
\]
This proves Theorem~\ref{thm:main}.

\begin{thebibliography}{9}
\bibitem{bonami} Aline Bonami, \emph{\'Etude des coefficients de Fourier des fonctions de $L^p(G)$}, Annales de l'Institut Fourier \textbf{20} (1970), no.~2, 335--402, \href{https://aif.centre-mersenne.org/item/AIF_1970__20_2_335_0/}{doi:10.5802/aif.357}.
\bibitem{odonnell} Ryan O'Donnell, \emph{Analysis of Boolean Functions}, Cambridge University Press, 2014, Theorem~9.21; \href{https://www.cs.cmu.edu/~odonnell/papers/Analysis-of-Boolean-Functions-by-Ryan-ODonnell.pdf}{author's complete manuscript}.
\end{thebibliography}
\end{document}
